\section{Planned Tables and Figures}
\label{sup:results}
\textcolor{red}{\textbf{Apart from the E1 rows marked in Table~\ref{tab:h6}, every table and figure below is a placeholder.}}

\todofig{repetition}{Cross-view correlation of aligned depth errors at the same reference surface point, by appearance class (textureless, specular, textured) and by frame separation along the strip, with the thresholds of H2 and the chosen block size $B$.}

\todofig{blocksweep}{Detection power, false-alarm rate and untestable fraction against the block size $B\in\{1,3,5,9\}$ on UseGeo and on MatrixCity with errors repeated across neighbouring frames (H2, H6).}

\begin{table}[h]
\caption{\textbf{[TODO]} Detection power (\%) at 5\% per-view false alarm on specular and textureless regional failures, natural and injected (H3).}
\label{tab:h3}
\centering\small\setlength{\tabcolsep}{3pt}
\begin{tabular}{lcccc}
\toprule
Method & Spec.\ nat. & Spec.\ inj. & Textl.\ nat. & Textl.\ inj. \\
\midrule
Aligned RMS            & \todocell & \todocell & \todocell & \todocell \\
Depth confidence       & \todocell & \todocell & \todocell & \todocell \\
Robust depth loss      & \todocell & \todocell & \todocell & \todocell \\
SpotLessSplats mask    & \todocell & \todocell & \todocell & \todocell \\
RobustNeRF mask        & \todocell & \todocell & \todocell & \todocell \\
Gate                   & \todocell & \todocell & \todocell & \todocell \\
Oracle                 & \todocell & \todocell & \todocell & \todocell \\
\bottomrule
\end{tabular}
\end{table}

\begin{table}[h]
\caption{\textbf{[TODO]} Depth-supervised splatting with and without the gate: held-out laser-depth AbsRel (lower is better) and surface F-score (higher is better) on failure and clean scenes; MatrixCity with injected faults is reported separately (H4).}
\label{tab:h4}
\centering\small\setlength{\tabcolsep}{3pt}
\begin{tabular}{lcccc}
\toprule
 & \multicolumn{2}{c}{UseGeo} & \multicolumn{2}{c}{ETH3D} \\
Method & AbsRel & F & AbsRel & F \\
\midrule
No gate                & \todocell & \todocell & \todocell & \todocell \\
Huber / GM depth loss  & \todocell & \todocell & \todocell & \todocell \\
RobustNeRF loss        & \todocell & \todocell & \todocell & \todocell \\
SpotLessSplats         & \todocell & \todocell & \todocell & \todocell \\
Gate                   & \todocell & \todocell & \todocell & \todocell \\
Oracle gate            & \todocell & \todocell & \todocell & \todocell \\
\bottomrule
\end{tabular}
\end{table}

\begin{figure}[h]
\centering
\input{\efigdir/fig_e1_runtime.tex}
\caption{E1 runtime: CPU wall-clock per stage, with the threads used and the seconds per gate pass (one
cavity test of every active view; pass counts derived from the saved loop rounds). A one-off load and
surfel build of the block took about 6 minutes and was cached.}
\label{fig:e1runtime}
\end{figure}

\todofig{calibration}{Histogram of $T_j/m_j$ on fault-free views, raw and with the leave-view-out scene temperature, the fitted $\hat\gamma$ per scene and the per-view heterogeneity $\omega$, and the false-alarm rate against $\alpha$ on test scenes, with and without deflation by $n_{\mathrm{eff}}/n$ or jackknife inflation (H5).}

\begin{table}[h]
\caption{\textbf{[TODO for the remaining rows]} Gate ablations: false-alarm rate and power (H6). Rows
with numbers are from E1 only: MatrixCity block~3, exact depth, three neighbouring frames carrying the
same $8\sigma$ error, noise $1.3\times$ the calibrated level, 10 seeds. FA is the share of clean views
excluded and power the share of faulty views excluded; not at matched false alarm. All other rows
remain placeholders.}
\label{tab:h6}
\centering\small
\begin{tabular}{lcc}
\toprule
Variant & FA (\%) & Power (\%) \\
\midrule
Full gate                         & 0.0 & 93.3 \\
Joint surface                     & \todocell & \todocell \\
Keep tested-view-only surfels     & \todocell & \todocell \\
Unwhitened (constant $\sigma$)    & 1.4 & 93.3 \\
Deflated by $n_{\mathrm{eff}}/n$  & \todocell & \todocell \\
FisherRF-style diagonal Fisher    & \todocell & \todocell \\
Bayes' Rays-style uncertainty     & \todocell & \todocell \\
Jackknife inflation               & \todocell & \todocell \\
No scene temperature              & 4.0 & 93.3 \\
No shared distance bias           & \todocell & \todocell \\
Fold-wise opacity warm-up         & \todocell & \todocell \\
Colour updates geometry (A9)      & \todocell & \todocell \\
Loop off (single pass)            & 6.1 & 93.3 \\
Drop all rejected per round       & \todocell & \todocell \\
Down-weight instead of drop       & \todocell & \todocell \\
Readmission                       & \todocell & \todocell \\
Single-view leave-out ($B=1$)     & 0.0 & 93.3 \\
Perturbed starting poses          & \todocell & \todocell \\
\bottomrule
\end{tabular}
\end{table}

\subsection{E1 Qualitative Checks}
\label{sup:e1qual}
The figures below are computed from the MatrixCity block already used in E1, rebuilt from the seed of one
grid scene ($8\sigma$ on 40\% of view 59); no new data and no training. Grey pixels are not testable,
because no surfel supported by other views covers them.

\begin{figure}[!htbp]
\centering
\input{\qfigdir/fig_e1_qual_views.tex}
\caption{E1 qualitative check of the test. Rows: a clean view far from the fault (view 69); the faulty
view; a clean neighbour on the same strip, swamped while the faulty view still shapes its cavity (view
58); and the same neighbour once the faulty view is dropped. Columns: true depth; the cavity's rendered
surfel depth; the raw residual after alignment; the residual standardised by $\sqrt{(\Sigma_r)_{uu}}$; and the
region-scan statistic $|z|$ per tile with the tiles above the first Holm threshold outlined. Flagged tiles:
0, 21, 16 and 0; largest $|z|$: 2.4, 80.1, 29.9 and 2.1.}
\label{fig:e1qualviews}
\end{figure}

\begin{figure}[!htbp]
\centering
\input{\qfigdir/fig_e1_qual_uncertainty.tex}
\caption{E1 residual uncertainty for a weakly supported view (view 56, 68\% of pixels testable) and a
strongly supported one (view 69, 100\%). Left: $\sqrt{(\Sigma_r)_{uu}}$, dominated by the depth noise,
which grows with range. Middle: its surface part $\sqrt{V_{-j}(u)}$, large where few other views support
the surfels. Right: the variance of each region-scan tile sum under the full within-surfel term
$UP^{(-j)}U^{\top}$ divided by its diagonal-only value. It is 1.0 on the strongly supported view; on the weakly
supported view the median is 1.54 and the largest 9.3. A diagonal-only test would understate those variances
by the same factors, which is the source of the false alarms that the full term removes.}
\label{fig:e1qualunc}
\end{figure}

\begin{figure}[!htbp]
\centering
\input{\qfigdir/fig_e1_qual_topdown.tex}
\caption{Top-down map of the surfels of MatrixCity block~3 (one 3~m cell per surfel voxel; highest surfel per
cell), coloured by height and by the normal-offset standard deviation after fusing all 185 views
(logarithmic; 30~m is the prior, where no view observes the surfel). Black: camera centres, an $8\times8$
grid of positions flown in four headings. Red: the views in Figures~\ref{fig:e1qualviews} and~\ref{fig:e1qualunc}.}
\label{fig:e1qualtop}
\end{figure}


\subsection{RGB Rendering Check}
\label{sup:rgb}
The gate uses depth only. As a check that the same block also supports a photometric surfel model, one
2D Gaussian splatting model~\citep{gs_2dgs} was trained with gsplat~\citep{gs_gsplat} on the 185 RGB training images of
MatrixCity block~3, all clean, at half resolution ($960\times540$). Training ran for 30\,000 iterations with an
L1 plus 0.2 D-SSIM loss and, after the first third, the normal-consistency and distortion terms of 2DGS. It
started from the 81\,152 E1 surfel centres, used the default gsplat densification, and was capped at
$1.5\times10^{6}$ Gaussians so that it fits an 8~GB GPU. It took 31~min on one RTX~2070 (peak GPU memory
3.2~GB; largest temperatures 66\,$^{\circ}$C GPU and 57\,$^{\circ}$C CPU). The 19 held-out test views of the
block lie between the training positions and point in headings that the training views do not use.
Table~\ref{tab:rgb} and Figure~\ref{fig:rgbtest} give PSNR and SSIM on them. The five test views at
$x=70$~m (frames 20--24) are the weakest, at 16.7--17.7~dB; in frame~22 the upper part of the image
lies beyond the area the training views cover and renders black (Figure~\ref{fig:rgbviews}, bottom row).

\begin{table}[!htbp]
\centering
\caption{2DGS RGB reconstruction of MatrixCity block~3: test-split mean $\pm$ standard deviation over
views, and the three training views of Figure~\ref{fig:e1qualviews}.}
\label{tab:rgb}
\input{\rfigdir/tab_rgb.tex}
\end{table}

\begin{figure}[!htbp]
\centering
\input{\rfigdir/fig_rgb_views.tex}
\caption{2DGS RGB renders against ground truth. Top three rows: the training views of
Figure~\ref{fig:e1qualviews} (their roles refer to the E1 depth test; the RGB images themselves are clean).
Bottom two rows: a well-reconstructed and the worst held-out test view. Right: absolute error, averaged over
the colour channels.}
\label{fig:rgbviews}
\end{figure}

\begin{figure}[!htbp]
\centering
\input{\rfigdir/fig_rgb_test.tex}
\caption{Per-view PSNR (\textbf{a}) and SSIM (\textbf{b}) of the 2DGS model on the 19 held-out views of
MatrixCity block~3; dashed: the mean.}
\label{fig:rgbtest}
\end{figure}
\FloatBarrier
